% =====================================================================
%  2bat-ud5-6.tex  —  SESSIÓ 6: Quants superaran el programa?
%  (sense preàmbul: s'inclou des de main.tex)
% =====================================================================
\horatitol{Sessió 6 — Quants superaran el programa?}%
{Activitat 2: modelitzar un grup real amb la binomial i traduir les probabilitats en decisions de planificació del servei.}

\subsection{Idea clau}

Apliquem la distribució binomial a un cas real (Activitat 2): sabem, per l'experiència
d'edicions anteriors, la probabilitat que un usuari superi un programa formatiu, i volem
\textbf{preveure el comportament del grup} per planificar el servei. Calcularem
probabilitats d'interès i, sobretot, les \textbf{traduirem en decisions}.

\begin{idea}
\textbf{Preguntes útils per a la planificació.} Amb $n$ usuaris i probabilitat d'èxit
$p$:
\begin{itemize}[leftmargin=1.4em, itemsep=2pt]
  \item \textbf{Nombre esperat:} $\mu=n\cdot p$ (quants superaran, de mitjana).
  \item \textbf{Exactament $k$:} $P(k)=\binom{n}{k}p^{k}(1-p)^{n-k}$.
  \item \textbf{Com a mínim $k$:} se sumen les probabilitats de $k$, $k+1$, \dots, $n$.
  \item \textbf{Com a màxim $k$:} se sumen les de $0$, $1$, \dots, $k$.
\end{itemize}
Cada probabilitat s'ha de traduir en una \textbf{conseqüència pràctica} (recursos,
places, continuïtat del programa).
\end{idea}

\noindent\textbf{El cas.} En edicions anteriors, la probabilitat que un usuari superi el
programa és $p=0,65$. Aquesta edició té un grup de $n=12$ usuaris.

\subsection{Exemples}

\begin{exemple}[nombre esperat i probabilitat exacta]
Amb $n=12$ i $p=0,65$, el nombre esperat d'usuaris que superaran el programa és
\[
\mu=12\cdot 0,65=7,8\approx 8.
\]
La probabilitat que en superin \textbf{exactament 8} és
\[
P(8)=\binom{12}{8}(0,65)^{8}(0,35)^{4}\approx 0,237=23,7\%.
\]
\end{exemple}

\begin{exemple}[com a mínim, per a la planificació]
Si el servei vol assegurar la fase següent només si superen el programa \textbf{com a
mínim 8} usuaris, cal la probabilitat $P(k\ge 8)$, que se suma de $8$ fins a $12$:
\[
P(k\ge 8)\approx 0,583=58,3\%.
\]
És a dir, hi ha una mica més d'un $50\%$ de possibilitats d'arribar al llindar: el servei
faria bé de preveure un pla alternatiu.
\end{exemple}

\subsection{Exercicis resolts}

\begin{exresolt}[exactament un nombre d'èxits]
En un grup de $10$ usuaris amb $p=0,8$ de completar un taller, quina és la probabilitat
que \textbf{exactament 9} el completin?
\solucio
$n=10$, $p=0,8$, $k=9$:
\[
P(9)=\binom{10}{9}(0,8)^{9}(0,2)^{1}=10\cdot 0,1342\cdot 0,2\approx 0,268=26,8\%.
\]
\resposta{$P(9)\approx 26,8\%$.}
\end{exresolt}

\begin{exresolt}[com a mínim un llindar]
Amb el mateix grup ($n=10$, $p=0,8$), quina és la probabilitat que en completin
\textbf{com a mínim 9} (és a dir, $9$ o $10$)?
\solucio
Se sumen $P(9)$ i $P(10)$:
\[
P(k\ge 9)=P(9)+P(10)\approx 0,268+0,107=0,376=37,6\%.
\]
\resposta{$P(k\ge 9)\approx 37,6\%$.}
\end{exresolt}

\subsection{Exercicis per fer}

\begin{exercicis}
\begin{exlist}
  \item Amb $n=15$ i $p=0,6$, calcula el nombre esperat d'usuaris que superaran el
        programa.
  \item Amb $n=6$ i $p=0,5$, calcula la probabilitat que en superin \textbf{exactament
        4}. (Pista: $\binom{6}{4}=15$.)
  \item Amb $n=8$ i $p=0,75$, calcula la probabilitat que \textbf{tots 8} completin el
        curs. (Pista: $k=8$.)
  \item Amb $n=20$ i $p=0,7$, calcula el nombre esperat d'èxits. Si el servei necessita
        almenys $14$ èxits, el valor esperat hi arriba?
  \item Amb $n=5$ i $p=0,4$, calcula la probabilitat que en tingui èxit \textbf{com a
        mínim 1}. (Pista: és més fàcil fer $1-P(0)$.)
  \item \textbf{(Interpretació)} Al cas del grup de $12$ amb $p=0,65$, s'esperen uns $8$
        èxits. Si cada usuari que supera passa a una fase següent que costa recursos,
        quants llocs caldria reservar per a aquesta fase? Per què convé preveure amb el
        valor esperat i no amb el màxim possible?
\end{exlist}
\end{exercicis}

% ---------------------------------------------------------------------
\fullsolucions{Sessió 6}
\begin{solucions}
\begin{sollist}
  \item $\mu=15\cdot 0,6=9$ usuaris.
  \item $P(4)=\binom{6}{4}(0,5)^4(0,5)^2=15\cdot 0,0625\cdot 0,25\approx 0,234=23,4\%$.
  \item $P(8)=\binom{8}{8}(0,75)^8(0,25)^0=(0,75)^8\approx 0,100=10,0\%$.
  \item $\mu=20\cdot 0,7=14$ èxits esperats. Sí: $14\ge 14$, arriba justament al llindar
        (tot i que, per ser el valor mitjà, convé tenir un marge).
  \item $P(k\ge 1)=1-P(0)=1-\binom{5}{0}(0,4)^0(0,6)^5=1-0,0778\approx 0,922=92,2\%$.
  \item Resposta oberta. Idea: es reservarien uns $8$ llocs (el nombre esperat). Preveure
        amb el valor esperat evita infra- o sobredimensionar el servei; planificar per al
        màxim possible ($12$) malbarataria recursos, i per al mínim en deixaria molts
        sense atenció. El valor esperat és l'estimació més raonable.
\end{sollist}
\end{solucions}
